\subsection{A uniform upper bound for two growing kernels}
\label{sec:growing-kernel-upper}

In this subsection the equilibrium is the original law $\pi_M$.
Let $\gamma_{\mathrm{prod}}$ and $\gamma_{\mathrm{add}}$ denote the
gaps obtained by using the shared multipliers $c(i,j)=ij$ and
$c(i,j)=i+j$, respectively.
Thus the product model has pair-merger rate $2ij/n$ and per-link
fragmentation rate $\beta ij$; its fragmentation rule and equilibrium
differ from those of the classical Ewens model above.

\begin{proposition}\label{prop:growing-kernel-upper}
For every $n\ge2$ and $\beta>0$,
\begin{align}
 \gamma_M
 &\le\gamma_{\mathrm{prod}}
       \le 2\beta+\frac{2(n-1)}n,\label{eq:product-kernel-bounds}\\
 2\gamma_M
 &\le\gamma_{\mathrm{add}}
       \le 4\beta+\frac{4(n-1)}n.\label{eq:additive-kernel-bounds}
\end{align}
In particular the growing kernels have gap bounded above by a
constant times $1+\beta$, uniformly in $n$.
\end{proposition}

\begin{proof}
The lower bounds follow from $ij\ge1$ and $i+j\ge2$ by
Dirichlet-form comparison.  To prove the upper bounds, write
$X=m_1$ for the number of monomers, $K=\sum_jm_j$, and
$\theta=\beta n$.
For the shared-product generator $L_{\mathrm{prod}}$, direct counting
gives
\begin{equation}\label{eq:monomer-product-drift}
 L_{\mathrm{prod}}X
   =2\beta(n-K)-\frac{2(n-1)}nX.
\end{equation}
Indeed a $j$-cluster with $j\ge3$ creates monomers at total rate
$2\beta(j-1)$.
A dimer splits at rate $\beta$ and creates two monomers, giving the
same contribution $2\beta(j-1)$.
Merging a monomer with every possible cluster, including another
monomer with its double loss, contributes $-2(n-1)X/n$.

We claim
\begin{equation}\label{eq:monomer-count-covariance}
 \operatorname{Cov}_{\pi_M}(X,K)\le\operatorname{Var}_{\pi_M}(X).
\end{equation}
Condition on $X=s$.  The remaining mass is $N=n-s$ and has no parts
of size one.  If $N=0$ its cluster count is zero.
The value $N=1$ is impossible under this conditioning.
For $N\ge2$, the unnormalized probabilities of its cluster count
$J_N=j$ are
\begin{equation}\label{eq:nonmonomer-count-weights}
 a_N(j)=\frac{\theta^j}{j!}\binom{N-j-1}{j-1},
 \qquad 1\le j\le\lfloor N/2\rfloor.
\end{equation}
To obtain this formula, sum $1/\prod_{r\ge2}m_r!$ over nonmonomer
partitions with $j$ parts.  Multiplying by $j!$ counts ordered
compositions of $N$ into $j$ parts each at least two, of which there
are $\binom{N-j-1}{j-1}$.

On the common support for $N,N+1$,
\[
 \frac{a_{N+1}(j)}{a_N(j)}
   =\frac{N-j}{N-2j+1},
\]
which is increasing in $j$: its derivative, viewed as a real
function, is $(N-1)/(N-2j+1)^2>0$.
Any newly available support point for $N+1$ is at the upper end.
Thus $J_{N+1}$ stochastically dominates $J_N$.
For completeness, an increasing likelihood ratio gives this
dominance by the nonnegative covariance of any increasing test
function with the likelihood ratio under the first distribution;
mass at a newly available largest point only increases such
expectations.
Also $J_0=0$ is stochastically below every $J_N$ with $N\ge2$.
It follows that the function
\[
 s\longmapsto\mathbb E[K-X\mid X=s]
       =\mathbb E J_{n-s}
\]
is decreasing on the feasible conditioning values
$s\in\{0,\ldots,n-2,n\}$.
The missing value $n-1$ causes no difficulty.
For a real random variable and a decreasing function of that
variable, covariance is nonpositive, as follows by taking an
independent copy and expanding the product of the two differences.
Consequently
\[
 \operatorname{Cov}(X,K)
 =\operatorname{Var}(X)
   +\operatorname{Cov}\bigl(X,\mathbb E[K-X\mid X]\bigr)
 \le\operatorname{Var}(X),
\]
proving \eqref{eq:monomer-count-covariance}.

The monomer count is nonconstant under $\pi_M$: both the all-monomer
and single-cluster states have positive mass.
Using it in the Rayleigh quotient and applying
\eqref{eq:monomer-product-drift} gives
\[
 \gamma_{\mathrm{prod}}
 \le\frac{-\operatorname{Cov}(X,L_{\mathrm{prod}}X)}
             {\operatorname{Var}(X)}
 =\frac{2(n-1)}n+
       2\beta\frac{\operatorname{Cov}(X,K)}{\operatorname{Var}(X)}
 \le\frac{2(n-1)}n+2\beta.
\]
Finally $i+j\le2ij$ for positive integers $i,j$.
The two chains have the same stationary law, so
$\mathcal E_{\mathrm{add}}\le2\mathcal E_{\mathrm{prod}}$ and
$\gamma_{\mathrm{add}}\le2\gamma_{\mathrm{prod}}$.
This proves the remaining upper bound.
\end{proof}

These estimates constrain the relaxation scale without reducing the
growing-kernel dynamics to a count chain.
At $\beta=b/n$ with fixed $b>0$, the comparison lower bounds have
order $n^{-1/2}$ while the upper bounds have constant order.
A lower bound of constant order in that regime is not established
by Proposition~\ref{prop:growing-kernel-upper}.
